% Chapter 6: Fleet Management
% Contains: Certificate Model, Workload Model, Clock Synchronization

\section{Certificate Model}
\label{sec:certificates}
% ============================================================

Certificates are JWTs (RFC 7519) signed with Ed25519 (\texttt{EdDSA} algorithm). A certificate attests that a DeviceId is valid for a time window.

\textbf{JWT Claims:}

\begin{lstlisting}[language=Rust]
struct CertClaims {
    sub: String,         // Subject DeviceId (base32)
    iss: String,         // Issuer DeviceId (base32)
    public_key: [u8; 32], // Subject's Ed25519 public key (base64)
    iat: i64,            // Issued-at timestamp (Unix seconds)
    nbf: i64,            // Not-before timestamp
    exp: i64,            // Expiration timestamp
}
\end{lstlisting}

\textbf{Self-signed root:} A root CA certificate has \texttt{sub == iss} and is verified against its own public key.

\textbf{Certificate chain validation:}

\begin{lstlisting}[language=Rust]
const MAX_CHAIN_DEPTH: usize = 10;

fn validate_cert(jwt: &str, validator: &CertValidator) -> Result<CertClaims> {
    let claims = decode_jwt_unsafe(jwt)?;  // Parse without signature check
    let issuer_pubkey = validator.get_pubkey(&claims.issuer_id()?)?;
    decode_and_verify_jwt(jwt, &issuer_pubkey)?;  // Verify signature
    validator.check_expiry(&claims)?;  // Check nbf <= now <= exp
    validator.verify_chain_to_root(&claims.issuer_id()?, 0)?;  // Recurse
    Ok(claims)
}
\end{lstlisting}

The validator caches intermediate certificates keyed by DeviceId. Chain validation recurses until reaching the trusted root or exceeding \texttt{MAX\_CHAIN\_DEPTH}.

\textbf{Delegated CA hierarchy:}

\begin{Verbatim}[frame=single,fontsize=\small]
Fleet Root CA (self-signed)
|-> Device Certificate (gateway-1)
|   roles: [gateway, workload_spawner, device_issuer]
|   subjects: [avena.>, app.>]
|   relay: full
|
|-> Device Certificate (tractor-1)
|   roles: [mobile, workload_spawner]
|   subjects: [avena.control.>, app.telemetry.>]
|   relay: subjects([avena.gossip.>, app.telemetry.>])
|
+-> Device Certificate (sensor-1)
    roles: [sensor]
    subjects: [app.telemetry.field1.>]
    relay: none
\end{Verbatim}

\textbf{Capability attenuation:} Issued certificates have capabilities $\subseteq$ issuer's capabilities.

\textbf{Certificate lifetime:} Months (configurable per fleet). Short relative to network connectivity bounds. Gossip-propagated revocation lists for emergencies.

\textbf{Certificate bootstrapping:} Initial device provisioning occurs out-of-band (e.g., USB, QR code, manufacturer pre-provisioning). Devices with the \texttt{device\_issuer} role may subsequently sign certificates for new devices, subject to capability attenuation constraints.

\textbf{Chain validation at handshake:} Each avenad instance is configured at startup with its trusted root CA certificate. During peer handshake, the certificate is included in the \texttt{HandshakeMessage}. Validation verifies: signature against issuer's public key, expiry bounds, and chain terminates at configured root.

\textbf{Revocation during partition:} Certificate revocation lists propagate via gossip. A revoked device can continue operating within its partition until either its certificate expires or the partition heals. Certificate lifetimes (months, not years) bound this exposure window.

% ============================================================
\section{Workload Model}
\label{sec:workloads}
% ============================================================

\begin{lstlisting}[language=Rust]
struct WorkloadSpec {
    id: WorkloadId,  // (issuer_id, name)
    target: WorkloadTarget,
    overrides: Option<WorkloadId>,

    // Container spec
    image: String,
    runtime: ContainerRuntime,
    env: HashMap<String, String>,
    mounts: Vec<MountSpec>,
    ports: Vec<PortSpec>,

    // Network connectivity mode
    network_mode: NetworkMode,

    // Authorization (scope for runtime interests)
    permitted_publish: Vec<SubjectPattern>,
    permitted_subscribe: Vec<SubjectPattern>,
    
    // Retention
    ttl: Duration,  // Default 6-12 months
    
    // CRDT metadata
    hlc: HybridTimestamp,
    signature: Ed25519Signature,
    certificate_chain: Vec<Certificate>,
}

enum WorkloadTarget {
    Device(DeviceId),
    Selector(DeviceSelector),
}

enum NetworkMode {
    Messaging,    // Pub/sub API only via host's NATS endpoint
    Addressable,  // Overlay IP with dedicated network sidecar
}
\end{lstlisting}

\textbf{Targeting:} Workloads target specific devices or match via selectors (roles, labels). Explicit targets can override selector-based workloads via the \texttt{overrides} field.

\textbf{Override semantics:} When workload $w'$ with $\text{target}(w') = \text{id}(d)$ specifies $\text{overrides}(w') = \text{id}(w)$, device $d$ receives $w'$ instead of $w$. This allows replacing a generic role-based workload with a specialized implementation on specific devices.

\textbf{Network mode:} Workloads declare their network connectivity requirements via \texttt{network\_mode}:

\begin{description}
    \item[Messaging (default):] Workload communicates exclusively through the pub/sub API. The workload connects to the host avenad's local NATS endpoint and publishes/subscribes to subjects. No overlay IP address is assigned. This mode is simpler, lighter, and sufficient for most application workloads.

    \item[Addressable:] Workload requires direct IP connectivity in the overlay network. The host avenad spawns an \texttt{avena-sidecar} container alongside the workload. The sidecar runs a userspace Wireguard endpoint with the workload's derived identity, providing an overlay IP address. Remote nodes can establish direct tunnels to the workload. The host routes Wireguard UDP packets to the sidecar but cannot decrypt workload traffic (isolation preserved via end-to-end encryption).
\end{description}

The Addressable mode is appropriate for workloads requiring raw IP connectivity (e.g., RTSP streams, custom protocols) rather than pub/sub messaging. Most workloads should use Messaging mode.

\textbf{TTL and garbage collection:} Specs expire after TTL. Garbage collection events broadcast via gossip. Tombstones have the same TTL (6--12 months).

% ============================================================
\section{Clock Synchronization}
% ============================================================

Avena encourages clock synchronization through multiple mechanisms:
\begin{enumerate}
    \item \textbf{NTP/PTP:} Devices should use NTP or PTP when network connectivity permits
    \item \textbf{GPS time:} Devices equipped with GPS receivers can derive precise time and may serve as NTP/PTP sources for nearby devices
    \item \textbf{Time service discovery:} The availability and location of time services propagates via gossip (under \texttt{avena.gossip.time.>})
    \item \textbf{HLC estimation:} When a device's wall clock diverges significantly from received HLC timestamps, the HLC algorithm naturally corrects by adopting the higher wall time
\end{enumerate}

\begin{lstlisting}[language=Rust]
struct TimeServiceInfo {
    device_id: DeviceId,
    service_type: TimeServiceType,  // NTP, PTP
    endpoint: SocketAddr,
    source: TimeSource,  // GPS, Network, Manual
    stratum: u8,
}
\end{lstlisting}

\subsection{Time Quarantine on Rejoin}

When a device rejoins after extended disconnection, clock drift may cause ordering anomalies. Avena implements time quarantine to prevent corrupted timestamps from propagating.

\textbf{On first peer connection after boot:}
\begin{enumerate}
    \item Compute $\Delta = |\text{local\_hlc.wall} - \text{peer\_hlc.wall}|$
    \item Apply threshold-based action:
\end{enumerate}

\begin{tabular}{@{}lll@{}}
\toprule
\textbf{Delta} & \textbf{Action} & \textbf{Rationale} \\
\midrule
$< 1$ minute & Normal operation & Drift within tolerance \\
$1$ min -- $1$ hour & Accept peer HLC, log warning & Minor drift, self-correcting \\
$> 1$ hour & \textbf{Quarantine} & Requires time sync \\
\bottomrule
\end{tabular}

\textbf{Quarantine behavior:}
\begin{itemize}[nosep]
    \item Device can \emph{receive} gossip state and application messages
    \item Device \emph{cannot} publish or create interests until time sync completes
    \item Exit quarantine via: GPS lock, NTP sync, or manual operator override
\end{itemize}

\begin{lstlisting}[language=Rust]
enum TimeStatus {
    Synchronized,
    Quarantined { delta_ms: u64 },
}
\end{lstlisting}

Thresholds are configurable per-fleet. V2X deployments may use tighter bounds (5 minutes); agricultural deployments may tolerate longer drift (4 hours).

